\chapter{Podsumowanie}\label{ch:podsumowanie}

\section{Wnioski}\label{sec:wnioski}

W podsumowaniu zwięźle przedstaw, w jakim stopniu zrealizowano cel określony we wstępie
(rozdział~\ref{ch:wstep}). Wnioski warto zapisać w formie wyliczenia:
\begin{enumerate}
    \item wniosek pierwszy, poparty wynikami z rozdziału~\ref{ch:wyniki},
    \item wniosek drugi, odnoszący się do przyjętej metodyki (rozdział~\ref{ch:metodyka}),
    \item wniosek trzeci, dotyczący ograniczeń zastosowanego podejścia.
\end{enumerate}

\section{Kierunki dalszych prac}\label{sec:dalsze-prace}

Wskaż możliwe rozszerzenia i usprawnienia, np.:
\begin{itemize}
    \item rozszerzenie badań o dodatkowe przypadki testowe,
    \item automatyzacja powtarzalnych czynności,
    \item porównanie z alternatywnymi rozwiązaniami.
\end{itemize}

\TemplateNote{Załączniki (np. dodatkową dokumentację, zrzuty ekranu, całe dokumenty PDF) umieszczaj w katalogu
\file{appendices/} i dołączaj w \file{main.tex}; po zmianie przełącznika \latexcmd{showappendicesfalse} w
\file{config/metadata.tex} załączniki znikają z dokumentu.}
